% !TeX spellcheck = da_DK
\input{../common/preamble.tex}

\title{Formandsguide}
\date{\formatdate{20}{4}{2026}}
\author{Casper Freksen\\
\small med tilføjelser af Oskar Haarklou Veileborg\\
\small \& Anders Bruun Severinsen}

\begin{document}

\maketitle

\section{Introduktion}
\label{sec:introduktion}

Dette er en guide omkring formandsarbejdet i \fredagscafeen.

Som formand er man foreningens billede udadtil, og har det overordnede ansvar for
at foreningen fungerer. 
Dertil kan der komme invitationer til diverse møder med instituttet, man bør deltage i.

Det er også vigtigt man har et overblik over, hvor de andre i
bestyrelsen er i deres opgaver. Det er ikke meningen, man skal ånde
dem i nakken, men hvis en falder bagud, er det godt at vide det og
tage en snak med dem, og sørge for at alle opgaver bliver udført eller fulgt op på.

\section{Bestyrelsesmøde}
\label{sec:bestyrelsesmode}

Hver måned afholdes der bestyrelsesmøde i \fredagscafeen. Her kan
man sende en afstemning ud et par uger før mødet bør holdes, hvor de forskellige medlemmer kan
vælge de datoer, de kan deltage. Send gerne en påmindelse en dag før afstemningen
lukkes. Endeligt bør der sendes en påmindelse på dagen for mødet.

Der tages referat til bestyrelsesmøderne af en referent. Se
evt. referaterne på GitLab repositoriet.

Det er en god ide at kigge tidligere referater igennem, for at huske
hvad der normalt tages op på mødet året før, og for at læse evt. eval
af tidligere afholdte events, så det kan tages til efterretning.

Dagsordenen bygges op mellem mødeindkaldelse og mødet. Hvis folk er
villige til at bruge GitHub Projects, kan de bidrage der. Ellers kan man sende
en mail med forslag til punkter.

Først laves der mad, så det kan være en god ide at sætte starttidspunktet for
mødes sådant at madlavning også indgår. 

Der er drikkevarer for ca. 50 kr. per person. 
Det ligger op til, at man deler med de andre, så man kan få smagt flere forskellige øl.

% Samtidig køres det store anlæg ud, og det skylles med lilla. 

Når mødet starter vælges dirigent og referent (statistisk set bliver
formanden valgt til mindst en af delene). Så behandles der punkter
ifølge dagsordenen.

\section{Generalforsamling}
\label{sec:generalforsamling}

Senest en måned inden den finder sted, skal der indkaldes til generalforsamling.

Generalforsamlingen er der, hvor du blev valgt. Da du gerne vil
genvælges, eller i det mindste gerne vil have en efterfølger, skal der
indkaldes til generalforsamlingen. Generalforsamlingen plejer at ligge
i starten af marts, så send derfor indkaldelsen i slutningen af januar.

Indkaldelsen skal indeholde en dagsorden, der f.eks. kunne se således ud:
\begin{itemize}
    \item Valg af dirigent og referent
    \item Beretning fra den afgående bestyrelse
    \item Beretning om foreningens regnskab
    \item Behandling af indkomne forslag
    \item Evaluering af foreningens arbejde
    \item Valg af formand
    \item Valg af kasserer
    \item Valg af resterende bestyrelse
    \item Valg af revisor
    \item Farvel til den gamle bestyrelse
    \item Eventuelt
\end{itemize}

Beretningen fra den afgående bestyrelse plejer formanden at tage sig
af. Forbered derfor et lille slideshow med billeder fra årets løb.

Der gives en øl til stemmeberettigede.

% Ift. til PR, så skal der gerne laves en plakat inden GF. Ifølge
% traditionen skal der være en generalhat på plakaten.

De afgående bestyrelsesmedlemmer skal gives en gave. 
Det har tidligere været en god flaske spiritus, 
eksempelvis whisky, rom, gin eller Ultrapres,
men spørg hver enkelt, hvad de har lyst til.
Prisen af gaven har tidligere været beregnet som 
en funktion $g$ over antal år i bestyrelsen $x$.
$$g(x) = 200 + 100x$$

GF kan være en fin mulighed for at tage et fællesbillede af den nye
bestyrelse.

\subsection{Efter generalforsamlingen}

Efter generalforsamlingen er der et par ting der skal gøres. De kan med fordel gøres
til første/overdragelses bestyrelsesmøde.

\begin{enumerate}
    \item[!] Bed de afgående bestyrelsesmedlemmer om at medbringe deres barnøgler.
    \item[!] Bed de nye bestyrelsesmedlemmer om at medbringe pas eller kørekort.
    \item[!] Alle forventes at medbringe studiekort, så der kan søges om den nødvendige adgang.
\end{enumerate}

\begin{itemize}
    \item De afgående bestyrelsmedlemmer afleverer deres nøgler tilbage, og de nye får
    udleveret deres. Dette noteres i oversigten over nøgler.
    \item Der skal skrives en mail til AU Card Access (\email{cardaccess@cs.au.dk}), om
    at få udvidet studiekort adgang til Institut for Datalogi, P-kælderen, ``Rummet ved 
    siden af rummet'' og opvaskerummet.
    I mailene skal der indgå en liste over de nye bestyrelsesmedlemmer's
    navne, studienumre, kortnumre og emails.\\
    Gæstekort skal også søges gennem bygningsdrift (\email{st.byg@au.dk}).
    \item Anmeldelse om ændring af sammensætningen af direktion og bestyrelse i selskaber, foreninger m.v., 
    der har alkoholbevilling udfyldes og sendes til Østjyllands Politi via
    Digital Post eller direkte til \email{ojyl-bevilling@politi.dk}.
    \item Ring til Danske Bank Foreninger \url{https://danskebank.dk/foreninger/find-hjaelp}, 
    og fortæl dem at du ønsker at meddele om ændring i bestyrelsen, og evt. i adgange til konti.
    De sender dig så en mail omkring hvilke dokumenter de har brug for, og hvordan du sender dem.\\
    Sidste gang bad de om følgende, og sendes på \href{denne side}{https://danskebank.dk/foreninger/send---hent-dokumenter}:
    \begin{itemize}
        \item Gældende vedtægter.
        \item Bestyrelsesoversigt – udfyldes online eller via en vedhæftede formular med den samlede bestyrelse, 
        inkl. navn, adresse, rolle og CVR nummer. 
        \item Legitimation på de tegningsberettigede (formand og kasserer) og dem der skal have fuldmagt (formand og kasserer). 
        Fra hver skal vi bede om billede af \textbf{sundhedskort}, sammen med gyldigt \textbf{kørekort eller pas}. 
        Bemærk, kopi fra apps accepteres ikke. 
        \item Hjælpeskema (vedhæftet) udfyldes med kontaktoplysninger og ønsker til aftalen samt hvem der skal slettes. 
    \end{itemize}
\end{itemize}

\subsection{Ny Formand}
\label{sec:ny-formand}

Dette skal kun gøres, hvis der er valgt en ny formand.
\begin{itemize}
    \item Tilføj den nye formand til virksomhedsregisteret på \url{virk.dk}. Husk at tildele
    alle de nødvendige adgangsrettigheder. Den tidligere formand kan med fordel stadig være koplet på,
    mens den tidligere tidligere formand kan fjernes.
    \item Husk at tilføje den nye formand til digital post, også på \url{virk.dk}.
    \item Gennemgang af og nøgle overleveres til pengeskabet i ``Rummet ved siden af rummet''.
    \item Aftale om lån af lokaler til afholdelse af fester/fredagsbarer skal underskrives og
    sendes til underskrift af Bygningschef for Administrationscenter Nat-Tech (\email{bent.lorenzen@au.dk}).
    \item Registrering af formandsskab for fredags bar/festforening på NT opdateres, underskrives
    af relevante parter og sendes til relevante interessenter.
    \item Registrering af forening/festforening på Aarhus Universitet opdateres, underskrives og sendes 
    til Nat-Tech Bygningsservice (\email{ntbyg.arrangementer@au.dk}).
    \item Ansøgning om godkendelse af bestyrer for virksomhed der har/anmoder om alkoholbevilling 
    til drift af restaurations- og hotelvirksomhed udfyldes og sendes til Østjyllands Politi via
    Digital Post eller direkte til \email{ojyl-bevilling@politi.dk}.
\end{itemize}

\subsection{Ny Kasserer}
\label{sec:ny-kasserer}

Dette skal kun gøres, hvis der er valgt en ny kasserer.
\begin{itemize}
    \item Tilføj den nye kasserer til digital post og fjern den tidligere kasserer, på \url{virk.dk}:
    \begin{enumerate}
        \item Log ind på \url{virk.dk} med MitID.
        \item Gå til ``Brugeradministration'' og vælg ``Gå til Digital Post for at se om du har adgang''.
        \item Vælg ``Rettigheder og adgang'' og derefter ``Tildel adgang''.
        \begin{enumerate}
            \item \textbf{Brugertype}: Vælg ``Medarbejder''.
            \item \textbf{Brugoplysninger}: Indtast den nye kasserers email og kaldenavn (Skriv deres fulde navn).
            \item \textbf{Adgangsrettigheder}: Vælg ``Basisadgang'' og ``Skriveadgang'' under ``Rettigheder i Digital Post''.
        \end{enumerate}
    \end{enumerate}
    \item Gennemgang af og nøgle overleveres til pengeskabet i ``Rummet ved siden af rummet''. 
    Dette kan med foredel også gøres af den tidligere kasserer.
\end{itemize}

\subsection{Porter Boarding}
\label{sec:porter-boarding}

`Porter Boarding' er navnet på et møde mellem den nuværende/nye
og ældre formænd, hvor der byttes viden for portere (eller andre mørke øl).
Du kan f.eks. invitere de ældre formænd til Porter Boarding en fredag i ugerne efter
generalforsamlingen.

Det kan anbefales at få nogle vise ord fra gamle formænd, og man har
mulighed for at komme med spørgsmål til, hvad man lige undrer sig
over. Derudover er det hyggeligt at snakke og få et par øl.

\section{Julegaver}
\label{sec:julegaver}

Vi plejer at give julegaver til drift og
rengøringen. Ift. drift har vi tidligere sponsoreret øl til deres
julefrokost (en fustage), men det var vidst lidt meget øl for dem, og
det er ikke blevet gjort de sidste par år. I stedet kan der blandes en
kurv med specialøl.

Til rengøring kunne der gives en æske god chokolade eller også nogle
blandede øl (hvis man ikke skal være så kønsdiskriminerende).

Endeligt kunne man overveje, om der er andre, vi gerne vil forkæle med
lidt ekstra op til jul.

\section{Andet}
\label{sec:andet}

Andre ting er ikke så fastlagt, men stadig en del af formandsarbejdet.

\begin{description}
    \item[Karantæner] Hvis en gæst eller bartender har opført sig imod barens adfærdsregler,
    kan det være nødvendigt at sætte vedkommende i karantæne. Dette
    skal besluttes af et flertal i bestyrelsen, og det er umiddelbart formandens
    ansvar at informere den pågældende om karantænen. Her kan man sende en mail
    fra hjemmesiden, så det er officielt.

    Nuværende karantæner kan findes på GitLab projektet under \texttt{/dokumenter} 
    \href{https://gitlab.au.dk/fredagscafeen/dokumenter/-/blob/main/Karantaener.md?ref_type=heads}{her}.

    \item[Førstehjælpskassen] Det er vigtigt at sørge for at
    førstehjælpskassen i baren indeholder de nødvendige ting. Dette
    inkluderer plaster, forbindinger, desinficerende middel
    og lignende. Tjek listen i kassen, og bestil nyt indhold
    når noget mangler. Tidlere har det været muligt at få nogle 
    af de mere specielle ting fra Proshop. Mere almindelige ting
    som plaster kan købes i Matas eller lignende.

    \item[Rusuge] Formanden kan være tovholder for introduktionen af fredagsbaren for russerne i rusugen.
    Efter intro til baren, gives en gratis sodavand til alle russerer og tutorer.

    \item[Barens fødselsdag] Det er tradition at vi giver en fustage pilsner på barens
    fødselsdag, som er den første fredag i september måned, som normalt er når russerne
    skal på Rustur.
\end{description}

\end{document}

%%% Local Variables:
%%% mode: latex
%%% TeX-master: t
%%% End:
